

%\documentclass[class=report , crop=false, multi={itemize, figure}, float=false]{standalone}
\documentclass{report}

%\usepackage{import} % Required for importing other .tex docs.  (import uses everything bw Begin and End Doc)
\usepackage{float} % Required for specifying the exact location of a figure or table
\usepackage{graphicx} % Required for including images
\usepackage{wrapfig}
\usepackage[pdftex,breaklinks,colorlinks=true,linkcolor=black,citecolor=blue,urlcolor=red,linktocpage=false,pagebackref=true,filecolor=magenta]{hyperref}%http://www.tug.org/applications/hyperref/manual.html#x1-100003.6
\usepackage{cite}
\usepackage[toc,title,page]{appendix}
\usepackage{pdfpages} % enables loading a pdf into the doc
\usepackage{makeidx}
\usepackage{glossaries} % must be after hyperref
\usepackage{blindtext}
\usepackage{enumitem}
%\usepackage{caption}

%\setlist[description]{leftmargin=\parindent,labelindent=\parindent}

%\renewcommand*{\bibname}{References} % renames the bibliography

\newcommand{\HRule}{\rule{\linewidth}{0.5mm}} % Command to make the lines in the title page

\graphicspath{{img/}{GIS_ChampionSection/img/}{awardsChapter/GIS_ChampionSection/img/}{brandPart/awardsChapter/GIS_ChampionSection/img/}{img/}{pairedProgSection/img/}{methodChapter/pairedProgSection/img/}{methodPart/methodChapter/pairedProgSection/img/}{documentationSection/img/}{methodChapter/documentationSection/img/}{methodPart/methodChapter/documentationSection/img/}{docStorageOrgSection/img/}{methodChapter/docStorageOrgSection/img/}{methodPart/methodChapter/docStorageOrgSection/img/}{QGisSection/img/}{toolsChapter/QGisSection/img/}{servicePart/toolsChapter/QGisSection/img/}{ESRISection/img/}{toolChapter/ESRISection/img/}{servicePart/toolChapter/ESRISection/img/}{../../../../source/}{../../source/}{servicePart/applicationsChapter/treasurerSection/img/}}

%\setlength\parindent{0pt} % eliminates indents


\usepackage[pdftex,breaklinks,colorlinks,citecolor=blue,urlcolor=red]{hyperref}

%\title{ % create title page
%\HRule % Horizontal Line added
%\\[.4cm] % space
%\begin{figure}[H] % included image
%\begin{center}	% centered horizontally
%	\includegraphics[scale=.45]{GIS_Logo_better.jpg}
%	\end{center}
%	\end{figure}
%	\Huge \bfseries TITLE GOES HERE % Title text
%\HRule \\[.4cm] % Horizontal Line added
%}  % closing brace for title
%
%\author{\Large Allegan County GIS \\\Large www.allegancounty.org/gis} % defines author

\begin{document}% Document Begins

%\ifstandalone
%%\frontmatter % turns off chapter numbering and uses roman numerals for page numbers
%\maketitle % creates title page and blank page after title page
%\tableofcontents % creates TOC and blank page
%\clearpage
%%\mainmatter % turns on chapter numbering, resets page numbering and uses arabic numerals for page numbers
%\fi

\subsection{hyperref Package}
\medskip 
\subsubsection{\Large usepackage}

Add the \textit{hyperref package} to the preamble\begin{large}\textbf{ last}\end{large}.	

\begin{verbatim}\usepackage[options]{hyperref}\end{verbatim}

\subsubsection{\Large Simple Use}
\begin{verbatim}Use \href{URL}{DESCRIPTION} to add a link with description
Example:
\href{https://www.latex-tutorial.com}{Website with tutorials}
produces:
\end{verbatim}
\href{https://www.latex-tutorial.com}{Website with tutorials}

	
\subsubsection{\Large Options}
Add optional arguments to the usepackage line:\\
Useful options:\begin{itemize}
\item \begin{large}\textbf{pdftex}\end{large}\\enables other options like breaklines
\item \begin{large}\textbf{breaklinks}\end{large}\\allow links to be broken across several lines\\eg. \url{https://lists.gnu.org/archive/html/emacs-orgmode/2013-06/msg00776.html}
\item \begin{large}\textbf{colorlinks}\end{large}\\Colors the text of links and anchors.(default is false)
\item \begin{large}\textbf{linkcolor}\end{large}\\Color for normal internal links(default is red).
\item \begin{large}\textbf{anchorcolor}\end{large}\\Color for anchor text.
\item \begin{large}\textbf{citecolor}\end{large}\\Color for bibliographic citations in text.
\item \begin{large}\textbf{urlcolor}\end{large}\\Color for linked URLs

\end{itemize}

\subsubsection{\Large Use with options}
\begin{verbatim}
\usepackage[breaklinks,colorlinks,citecolor=blue,urlcolor=green]{hyperref}
\end{verbatim}

\subsubsection{\Large Commands}
\begin{verbatim}
\href{URL}{text} — Makes text a link to URL.
\end{verbatim}
\href{../../../../../documentation/references/hyperrefmanual.pdf}{Manual for hyperref}
%\href{../hyperrefmanual.pdf}{Manual for hyperref}
%\href{c://AC/jalapeno/documentation/hyperrefmanual.pdf}{Manual for hyperref}

\begin{verbatim}


\hyperref[label]{text} 
  — Makes text a link to where
\ref{label} would point.
\hypertarget{name}{text} — Sets an anchor on textwith the label name. 
\hyperlink{name}{text} — Makes text a link that takes you to the anchor labeled name. 
Pair with \hypertarget. 
\phantomsection — Used in conjunction with
\addcontentsline to make the correct link in the Table of Contents.
\end{verbatim}				


\end{document}
